
\section{Limitations}

ARES relies on a small set of annotations in the human preference validation set (roughly 150-300 datapoints but more is better). 
These annotations often require an annotator familiar with the RAG system's domain application.
While these annotations can be easy to generate for general-domain applications, more specialized domains, such as law, medicine, and finance, may require annotators with specialized expertise.

The LLMs used in ARES benefit substantially from GPU-based hardware with substantial storage. 
In ARES, DeBERTa-v3-Large (304M) and FLAN-T5-XXL (11.3B) required GPUs with about 32GB of memory to run, taking several hours for fine-tuning and generation, respectively.
While commercial GPUs are widely available, they are not easily accessible to all NLP researchers and practitioners due to their costs.

Additionally, all of the datasets used in our evaluation of ARES are in English, a well-resourced language with abundant annotations.
Future work should explore how ARES can be employed in other languages by utilizing different LLMs for the ARES judge and the synthetic data generation.
This can help us better understand the strengths and weaknesses of the current ARES framework.
